\chapter{Общая следовая нормировка cross-carrier вложения}
\label{ch:v8-baryon-c0-common-trace-embedding-normalization-gate}

Предыдущий гейт свёл свободу карты $c_0$ к положительному множителю
$\kappa$. Проверим, фиксирует ли его требование общего следа. До построения
межсекторной стрелки минимальная общая алгебра является прямой суммой
\begin{equation}
 \mathcal A_{\rm sep}=\mathcal A_{\rm src}\oplus\mathcal A_{\rm aux}.
 \label{eq:v8-baryon-c0-separated-common-algebra}
\end{equation}

\section{Симплекс общих следов}

Даже после нормировки единицы положительный след на двух центральных
суммандах имеет семейство
\begin{equation}
 \tau_{p,q}(X\oplus Y)
 =p\,\tau_{\rm src}(X)+q\,\tau_{\rm aux}(Y),
 \qquad p,q>0,qquad p+q=1.
 \label{eq:v8-baryon-c0-direct-sum-trace-simplex}
\end{equation}
Пусть внутренние генераторы прямых $u,e$ нормированы в собственных следах:
\begin{equation}
 \tau_{\rm src}(u^*u)=1,
 \qquad
 \tau_{\rm aux}(e^*e)=1.
 \label{eq:v8-baryon-c0-summand-unit-norms}
\end{equation}
Тогда общая следовая метрика даёт
\begin{equation}
 G_{\rm src}(u,u)=p,
 \qquad
 G_{\rm aux}(e,e)=q.
 \label{eq:v8-baryon-c0-weighted-line-metrics}
\end{equation}

Для $\mathcal M_\kappa(u)=\kappa e$ условие pullback-изометрии равно
\begin{equation}
 \mathcal M_\kappa^*G_{\rm aux}\mathcal M_\kappa=G_{\rm src}
 \quad\Longleftrightarrow\quad
 q\kappa^2=p,
 \label{eq:v8-baryon-c0-trace-isometry-equation}
\end{equation}
то есть
\begin{equation}
 \kappa=\sqrt{\frac pq}.
 \label{eq:v8-baryon-c0-trace-weight-selector}
\end{equation}
След выбирает $\kappa$ только после выбора центрального отношения $p/q$.

\section{Два точных общих следа}

Два нормированных положительных следа дают разные изометрические карты:
\begin{equation}
 (p,q)=\left(\frac12,\frac12\right)\Longrightarrow\kappa=1,
 \qquad
 (p,q)=\left(\frac45,\frac15\right)\Longrightarrow\kappa=2.
 \label{eq:v8-baryon-c0-two-trace-witnesses}
\end{equation}
Оба сохраняют калибровочную инвариантность, положительность и нормировку
$\tau(1)=1$. Поэтому слово ``общий'' не устраняет центральную свободу.

Она видна непосредственно на центральных проекторах
\begin{equation}
 P_{\rm src}=1\oplus0,
 \qquad P_{\rm aux}=0\oplus1,
 \qquad
 \tau_{p,q}(P_{\rm src},P_{\rm aux})=(p,q).
 \label{eq:v8-baryon-c0-central-projector-weights}
\end{equation}

\section{Как свобода может исчезнуть}

Центральные веса перестают быть независимыми только после замены прямой
суммы настоящей linking-алгеброй с ненулевым бимодулем:
\begin{equation}
 \mathcal L=
 \begin{pmatrix}
 \mathcal A_{\rm src}&\mathcal E^*\\
 \mathcal E&\mathcal A_{\rm aux}
 \end{pmatrix},
 \qquad \mathcal E\ne0.
 \label{eq:v8-baryon-c0-required-linking-algebra}
\end{equation}
В простом полном матричном блоке единый нормированный trace действительно
связывает веса углов их размерностями. Но существующий барионный auxiliary
carrier ещё не включён в такой блок вместе с источником $r$.

\begin{theorem}[Direct-sum common-trace no-go]
На разделённой алгебре общий нормированный положительный след образует
одномерный симплекс центральных весов. Pullback-изометрия фиксирует
$\kappa=\sqrt{p/q}$, но отношение $p/q$ остаётся свободным. Поэтому общий
след без linking-бимодуля не выбирает $\kappa$ и $c_0$.
\end{theorem}

\begin{nogocriterion}[Запрет скрытого равенства центральных весов]
Назначение $p=q$ является дополнительным выбором, если не предъявлен
оператор, смешивающий два угла в одном простом следовом блоке. Равенство
размерностей прямых само по себе не связывает центральные состояния
прямой суммы.
\end{nogocriterion}

\begin{protocol}\begin{enumerate}
\item общий положительный нормированный trace: \textbf{существует};
\item размерность симплекса весов: \textbf{$1$};
\item условная изометрическая норма: \textbf{$\kappa=\sqrt{p/q}$};
\item $p/q$ выведено: \textbf{нет};
\item $\kappa$ и $c_0$ выбраны: \textbf{нет};
\item следующий гейт: \textbf{допуск off-diagonal linking bridge}.
\end{enumerate}\end{protocol}

Машинный сертификат сохранён в
\path{s2t_v8_baryon_c0_common_trace_embedding_normalization_gate_results.json}.